\section{Geometric Control Theory}
\begin{definizione}
	The system \ref{eq:2.2.1} is said to be controllable if for any two points $x_0$ and $x_f$ in M there exist and admissible controlo $u(t)$ defined on some time interval $[0,T]$ such that the system $\ref{eq:2.2.1}$ with initial condition $x_0$ reaches the point $x_f$ in time T
\end{definizione}
Generally it's not easy to prova that a system is contrabble , so we introduce a related property called \textbf{accessability} , whisch is often much easier to prove . We first need the notion of reachable set
\begin{definizione}
	Given $x_0 \in M$ we define $R(x_0,t)$ to be the set of all $x\in M$ for which there exists an admissinle controlo u such taht there is a trajectory of  \ref{eq:2.2.1} with $x(0)=x_0 \quad , \quad x(t))=x$. The reachable set form $x_0$ at time T is definied to be 
	\[R_T(x_0) = \bigcup_{0\leq t\leq T} R(x_0,t)\]
\end{definizione}
\begin{definizione}
	The accessibility algebra C of \ref{eq:2.2.1} is the smallest Lie Algebra of vector field onb M that constains the vector Fields f and $g_1 , \dots , g_m$ 
\end{definizione}
The Accessibility algebra is just the span of all possible Lie brackets of f and the $g_i$
\begin{definizione}
	We define the accessibility distribution C of \ref{eq:2.2.1} to be the distribution generated by the vector field in C , i.e. , $C(x)$ is the span of the vector field X in  C at z 
\end{definizione}
\begin{definizione}
The system \ref{eq:2.2.1} is said to be accessible from $p \in M $ if for evey $T > 0 $ , $R_T(p)$ contains a non empty open set
\end{definizione}
\begin{teo}{}{}
Consider the system \ref{eq:2.2.1} and assume that the vector field are $C^\infty$. If dim$C(x_0) = n$ ( the accessibility algebra span the tangent space to M at $x_0$) the for any $T>0$ the set $R_t(x_0)$ has a nonempry interior ; i.e. , the system ha the accessibility propery from $x_0$
\end{teo}
\begin{definizione}
A \textbf{system} or \textbf{machine} $\Sigma = (\mathcal{T}, X, \mathcal{U}, \phi)$ consists of:
\begin{itemize}
	\item A time set $\mathcal{T}$;
	\item A nonempty set $X$ called the \textbf{state space} of $\Sigma$;
	\item A nonempty set $\mathcal{U}$ called the \textbf{control-value} or \textbf{input-value space} of $\Sigma$; and
	\item A map $\phi : \mathcal{D}_\phi \to X$ called the \textbf{transition map} of $\Sigma$, which is defined on a subset $\mathcal{D}_\phi$ of
	\begin{equation}
		\{(\tau, \sigma, x, \omega) \mid \sigma, \tau \in \mathcal{T}, \sigma \le \tau, x \in X, \omega \in \mathcal{U}^{[\sigma, \tau)}\},
	\end{equation}
\end{itemize}
such that the following properties hold:

\medskip
\noindent \framebox{\textbf{nontriviality}} For each state $x \in X$, there is at least one pair $\sigma < \tau$ in $\mathcal{T}$ and some $\omega \in \mathcal{U}^{[\sigma, \tau)}$ such that $\omega$ \textbf{is admissible for} $x$, that is, so that $(\tau, \sigma, x, \omega) \in \mathcal{D}_\phi$;

\medskip
\noindent \framebox{\textbf{restriction}} If $\omega \in \mathcal{U}^{[\sigma, \mu)}$ is admissible for $x$, then for each $\tau \in [\sigma, \mu)$ the restriction $\omega_1 := \omega|_{[\sigma, \tau)}$ of $\omega$ to the subinterval $[\sigma, \tau)$ is also admissible for $x$ and the restriction $\omega_2 := \omega|_{[\tau, \mu)}$ is admissible for $\phi(\tau, \sigma, x, \omega_1)$;

\medskip
\noindent \framebox{\textbf{semigroup}} If $\sigma, \tau, \mu$ are any three elements of $\mathcal{T}$ so that $\sigma < \tau < \mu$, if $\omega_1 \in \mathcal{U}^{[\sigma, \tau)}$ and $\omega_2 \in \mathcal{U}^{[\tau, \mu)}$, and if $x$ is a state so that
\[ \phi(\tau, \sigma, x, \omega_1) = x_1 \quad \text{and} \quad \phi(\mu, \tau, x_1, \omega_2) = x_2, \]
then $\omega = \omega_1 \omega_2$ is also admissible for $x$ and $\phi(\mu, \sigma, x, \omega) = x_2$;

\medskip
\noindent \framebox{\textbf{identity}} For each $\sigma \in \mathcal{T}$ and each $x \in X$, the empty sequence $\diamond \in \mathcal{U}^{[\sigma, \sigma)}$ is admissible for $x$ and $\phi(\sigma, \sigma, x, \diamond) = x$. \hfill $\square$
\end{definizione}
\begin{definizione}
	The system $\Sigma$ is locally controllable in time T at the equilibrium state x if and only if for every neighborhood $\mathcal{V} $ of x there is another neighboord $\mathcal{W}$ of x , such that , for any pair of elements z , y in W , z can be controlled to y inside $\mathcal{V}$ in time T 
\end{definizione}
\begin{teo}{}{}
 Assume that $\Sigma$ is a continuos-time system of class $C^1$ . Then , 
 \begin{enumerate}
 	\item Let $\Gamma = (\xi, \omega)$ be a trajectory for $\Sigma$ on an interval $I = [\sigma , \tau]$ . Thean , a sufficient condition for $\Sigma$ to be locally controllable along $\xi$ is that $\sigma_* [\Gamma]$ be controllable on $[\sigma , \tau]$ 
 \item Ler $\Sigma$ be time-invariant , and assume that x is an equilibrium state . Pick ant $\varepsilon >0$ and any $u \in U$ such that $f(x,u) = 0$ . Then , a sufficient condition for $\Sigma$ to be locally controllable at in time $\epsilon$ is that the linearization of $\Sigma$ at (x,u) be controllable
  \end{enumerate}
\end{teo}
We now define the controllabilty of a linear sisystem 
\begin{teo}{}{}
	The n-dimensionale continuos and invariant time  linear system $\Sigma$ is controllable if and only if $rank R(A,B) = n  $ where 
	\[R(A,B) = [ B , AB,A^2B,\dots , A^{n-1}B]\]
\end{teo}

